
% ============================================================
% TABLE 09 — SCOPE AND LIMITS OF INVARIANCE
% Book I: Ontology of Viability
% Part I: Identity and Manifestation
% Chapter 2: Invariance
% ============================================================

\begingroup

\small
\renewcommand{\arraystretch}{1.25}

\setlength{\LTleft}{0pt}
\setlength{\LTright}{0pt}
\setlength{\LTpre}{1.0em}
\setlength{\LTpost}{1.0em}

\begin{longtable}{
    @{}
    >{\raggedright\arraybackslash}p{0.18\textwidth}
    >{\raggedright\arraybackslash}p{0.36\textwidth}
    >{\raggedright\arraybackslash}p{0.36\textwidth}
    @{}
}

\caption{The Scope and Limits of the Concept of Invariance}
\label{tab:scope-and-limits-of-invariance}
\\

\toprule
\textbf{Dimension} &
\textbf{What the Concept Establishes} &
\textbf{What It Does Not Establish}
\\
\midrule
\endfirsthead

\multicolumn{3}{@{}l@{}}{
    \textit{The Scope and Limits of the Concept of Invariance (continued)}
}
\\
\toprule
\textbf{Dimension} &
\textbf{What the Concept Establishes} &
\textbf{What It Does Not Establish}
\\
\midrule
\endhead

\midrule
\multicolumn{3}{r@{}}{
    \textit{Continued on the next page}
}
\\
\endfoot

\bottomrule
\endlastfoot

Subject of invariance &
An existent, or a determinate aspect of an existent,
may be considered with respect to invariance. &
The subject of invariance is not thereby reduced
to the particular determination under consideration.
\\

\addlinespace[0.5em]

Domain of change &
Invariance is assessed in relation to a specified
domain of changes or transformations. &
Invariance within one domain does not establish
invariance under every possible change
or transformation.
\\

\addlinespace[0.5em]

Respect of invariance &
A particular determination, feature, or relation
may remain unchanged under the specified
transformations. &
The invariance of one respect does not establish
the invariance of every aspect of the existent.
\\

\addlinespace[0.5em]

Identity &
Invariance may help describe how a determinate
aspect remains unchanged under specified
transformations. &
Invariance does not constitute identity, provide
its complete explanation, or establish that
identity is exhausted by invariant features.
\\

\addlinespace[0.5em]

Difference &
Different determinations may coexist with
the invariance of a particular determination. &
Difference does not necessarily negate invariance,
and invariance does not eliminate difference.
\\

\addlinespace[0.5em]

Change &
An existent may undergo change in some respects
while remaining invariant in another respect. &
Invariance does not mean the total absence
of change or the immutability of the existent
as a whole.
\\

\addlinespace[0.5em]

Continuity &
Invariance may contribute to the description
of certain patterns of continuity. &
The invariance of a particular feature does not,
by itself, establish the continuity of the existent
as a whole.
\\

\addlinespace[0.5em]

Structure &
Particular structural determinations or relations
may remain invariant under specified
transformations. &
The invariance of one structural relation does not
establish the invariance of the entire structure.
\\

\addlinespace[0.5em]

Knowledge &
Invariance may be investigated, recognized,
compared, and formally described through
epistemic operations. &
Epistemic recognition does not constitute
ontological invariance, and the absence
of recognition does not establish its absence.
\\

\addlinespace[0.5em]

Ontological significance &
The relational concept provides an initial
framework for describing determinate invariance. &
This framework does not settle whether invariance
has a deeper ontological status beyond the
relational account adopted in this edition.
\\

\end{longtable}

\endgroup

% ============================================================
% END OF TABLE 09
% ============================================================
